\section{线性化：把弯的暂时当成直的}
\label{sec:18-Synthesis/02-Recurring-Ideas/01}

同一个下午撞见两件事。一件是解方程：Newton 迭代从初值 \(1.2\) 出发三步收敛到八位有效数字，把初值改成 \(1.6\)，第三步跳到 \(-10^{4}\)，第四步溢出。另一件在隔壁桌：一台移动机器人的状态估计，走直线时轨迹漂亮，一进急转弯，位置估计偏出去半米，而滤波器报出的协方差反而比直线段更小——它比任何时候都更自信，也比任何时候都更错。

一个是求根，一个是传感器融合，症状一个是发散一个是过度自信，看不出有什么关系。但两段代码里都有同一步操作：在当前点把一个非线性函数换成它的切线（切平面），然后拿这条直线当作真函数来用。两处都在那一步失守。Newton 法把切线的零点当成函数的零点，初值落在拐点附近时切线几乎水平，零点被甩到很远；扩展 Kalman 滤波把非线性观测模型在当前估计处线性化，急转弯时真实映射的弯曲被那条切平面完全丢掉，而协方差的递推只用了 Jacobian，它记录的是那条直线的不确定性，不是真实映射的。

\subsection{一个动作}

这本书里出现次数最多的数学动作，不是求导也不是积分，是这一句：\emph{在一点附近，用一阶项代替原函数。}Taylor 展开把它写成一句可以清账的话，
\[
  f(x_0+h) \;=\; f(x_0) + f'(x_0)\,h + \tfrac{1}{2} f''(\xi)\,h^{2},
  \qquad \xi \text{ 介于 } x_0 \text{ 与 } x_0+h \text{ 之间},
\]
多元的版本把 \(f'(x_0)\) 换成 Jacobian 或梯度，把 \(f''(\xi)\) 换成 Hessian。式子右端的前两项就是线性化的产物，第三项是被扔掉的东西。所有关于线性化的讨论——它什么时候可信、什么时候骗人、邻域该取多宽——都是关于第三项的讨论。

\begin{center}
\begin{tikzpicture}[x=1cm,y=0.9cm,font=\small,>=stealth]
  % ---- 左图：曲线与切线，偏差随距离增长 ----
  \draw[black!35,->] (-0.2,0) -- (6.0,0) node[right,black!60,font=\footnotesize] {\(h\)};
  \draw[black!35,->] (0,0) -- (0,5.0);
  \draw[blue!65!black,very thick,domain=0:5.6,samples=90,smooth]
    plot (\x,{0.22*\x*\x - 0.5*\x + 1.6});
  \draw[red!70!black,thick,dashed,domain=0:5.6,samples=2]
    plot (\x,{1.48 + 0.38*(\x-2)});
  \draw[black!30,dotted] (2,0) -- (2,1.48);
  \fill[black] (2,1.48) circle (0.06);
  \node[anchor=north,font=\footnotesize,black!70] at (2,-0.05) {\(x_0\)};
  \draw[black!70,|-|] (3,1.86) -- (3,2.08);
  \node[anchor=west,font=\footnotesize,black!75] at (3.08,1.97) {小};
  \draw[black!70,|-|] (5,2.62) -- (5,4.60);
  \node[anchor=west,font=\footnotesize,black!75] at (5.08,3.6) {\(\tfrac{1}{2}f''h^{2}\)};
  \node[anchor=west,font=\footnotesize,blue!65!black] at (0.15,4.5) {真函数};
  \node[anchor=west,font=\footnotesize,red!70!black] at (0.15,3.95) {切线};

  % ---- 右图：曲率决定可信邻域宽度 ----
  \begin{scope}[shift={(7.1,0)}]
    \draw[black!35,->] (-0.2,0) -- (4.3,0);
    % 平缓
    \fill[green!16] (0.78,3.05) rectangle (3.22,4.55);
    \draw[blue!65!black,very thick,domain=0:4,samples=60,smooth]
      plot (\x,{0.10*(\x-2)*(\x-2) + 3.6});
    \draw[red!70!black,thick,dashed] (0,3.6) -- (4,3.6);
    \node[anchor=west,font=\footnotesize,black!75] at (0.05,4.85) {曲率小：可信邻域宽};
    % 陡峭
    \fill[green!16] (1.59,0.35) rectangle (2.41,1.85);
    \draw[blue!65!black,very thick,domain=0.6:3.4,samples=60,smooth]
      plot (\x,{0.90*(\x-2)*(\x-2) + 0.8});
    \draw[red!70!black,thick,dashed] (0,0.8) -- (4,0.8);
    \node[anchor=west,font=\footnotesize,black!75] at (0.05,2.15) {曲率大：可信邻域窄};
  \end{scope}
\end{tikzpicture}

\emph{左图：切线的误差按 \(h^{2}\) 增长，所以``离得近''这句话必须换算成具体的距离。右图：容许同样大的误差，曲率大的地方允许的步子小得多——绿色带子的宽度就是这一节要反复用到的``可信邻域''。}
\end{center}

\subsection{可信邻域有多宽}

把上面那张图读成数字。要求被丢掉的二阶项相对一阶项不超过 \(\varepsilon\)，就是要求
\[
  \tfrac{1}{2}\abs{f''}\,h^{2} \;\le\; \varepsilon\,\abs{f'}\,\abs{h}
  \qquad\Longleftrightarrow\qquad
  \abs{h} \;\le\; \frac{2\varepsilon\,\abs{f'}}{\abs{f''}} .
\]
右端这个量就是可信邻域的半径。它不是一个常数，而是随位置变化的——梯度越平、曲率越大的地方，它越小。多元情形把它读成 \(\norm{\nabla f}\) 与 Hessian 谱范数之比，结论一样：曲率大的方向，允许的步子短。

这个式子里有一件事值得盯住：分母是曲率，分子是斜率。所以线性化最先失效的地方不是曲率最大的地方，而是\emph{曲率相对斜率最大}的地方——也就是函数快要转平的地方。Newton 法在拐点附近发散，梯度下降在鞍点附近爬行，delta 方法在函数导数接近零时失准，说的都是分子塌了而不是分母涨了。\hyperref[sec:02-Calculus/05-Sequences-and-Series/06]{``Taylor 展开与余项：匹配局部变化信息''}把余项的各种写法列全了，实际使用中真正要估的从来只是这一个比值。

\subsection{同一个动作，八个名字}

导数这个概念本身就是线性化的定义，不是它的应用。\hyperref[sec:10-Matrix-Calculus/01-Derivative-Conventions/01]{``导数是最佳线性映射''}把这句话说到底：所谓可微，就是存在一个线性映射使得余项相对增量趋于零；导数是那个映射，不是那个数。一元情形下它是切线的斜率，多元情形下它是 Jacobian，\hyperref[sec:02-Calculus/06-Multivariable-Calculus/04]{``Jacobian 与多元链式法则''}说明复合函数的线性化就是各层线性化的矩阵乘积——反向传播算的正是这个乘积，只不过按从右往左的顺序结合以节省计算。

Newton 法把线性化用在方程的解上：不解 \(f(x)=0\)，解它的切线的零点，再从那里重来。\hyperref[sec:08-Numerical-Analysis/02-Root-Finding/02]{``不动点与 Newton 法''}给出它二次收敛的来源——每一步的误差被平方，正是因为丢掉的恰好是二阶项。同一个动作用在优化上，就是在梯度为零的方程上做 Newton 法，即\hyperref[sec:09-Convex-Optimization/06-Optimization-Algorithms/02]{``Newton 法用曲率校正方向''}：这时被线性化的是梯度，被保留的是 Hessian。

梯度下降只用一阶项决定方向，然后靠步长把自己限制在可信邻域内。\hyperref[sec:09-Convex-Optimization/06-Optimization-Algorithms/01]{``梯度下降真正难在步长''}那一篇的核心结论——步长上界由曲率的上界给出，收敛速度由曲率的最大最小之比给出——就是上面那个半径公式在优化语境里的复述。那个比值叫\hyperref[sec:08-Numerical-Analysis/01-Floating-Point-and-Error/03]{``条件数衡量问题本身的敏感性''}里的条件数，它同时管着数值求解的误差放大和梯度法的震荡幅度，因为两处问的是同一个问题：一阶信息在各个方向上的可信半径差多少倍。

扩展 Kalman 滤波把非线性的运动模型与观测模型在当前估计处线性化，然后套用线性情形的递推。\hyperref[sec:16-Control-and-Reinforcement-Learning/01-Optimal-Control/04]{``Kalman 滤波：用预测与更新的循环从噪声中估计状态''}讲清了线性 Gaussian 情形下这套递推为什么是精确的最优估计；一旦线性化进来，最优性就没有了，剩下的只是一个可以跑的近似。\hyperref[sec:16-Control-and-Reinforcement-Learning/01-Optimal-Control/06]{``DDP：用局部二次模型改进整条轨迹''}做的是同一件事的控制版本：沿着当前轨迹把动力学线性化、把代价二次化，解一个局部的线性二次问题，再走一小步。它必须配一个信赖域，理由与上面完全一样。

动力系统的稳定性判据是线性化最省力的一次使用。要判断一个不动点附近的轨迹是收还是发，不必解方程，只要把向量场在不动点处线性化，看 Jacobian 的谱落在虚轴哪一侧。\hyperref[sec:15-Differential-Equations-and-Dynamical-Systems/02-Stability-and-Dynamical-Systems/02]{``稳定性与 Lyapunov 函数''}给出这条判据成立的条件，也给出它失效的位置：谱恰好落在虚轴上时，一阶项决定不了任何事，非线性项说了算。

统计里的 delta 方法把估计量的方差通过一阶展开传过去：若 \(\hat\theta\) 渐近正态，则 \(g(\hat\theta)\) 的渐近方差是 \(g'(\theta)^{2}\) 乘以 \(\hat\theta\) 的渐近方差。\hyperref[sec:11-Mathematical-Statistics/04-Interval-Estimation/03]{``渐近区间与 Bayesian 可信区间''}用它构造区间。这里被线性化的是变换函数，可信邻域的宽度由估计量自己的抽样波动决定——样本量小，波动大，波动就可能超出 \(g\) 保持近似线性的范围。

到了深度学习，同一个动作出现在两个完全不同的位置。局部可解释方法在一个样本附近拟合一个线性模型，用它的系数解释黑箱的这次输出，\hyperref[sec:14-Deep-Learning/10-Interpretability-Mathematics/04]{``模型解释的数学形式化''}把这类方法的承诺与限度摆开讲。神经切核则把整个训练早期的网络当成关于参数的线性模型——参数变化小，函数就近似是参数的线性函数，训练动力学退化为核回归，\hyperref[sec:14-Deep-Learning/10-Interpretability-Mathematics/02]{``特征学习与核方法的差异''}正是从这条线性化能撑多久开始，讨论特征学习何时才真正发生。

\subsection{失效是同一种失效}

上面八处的报错方式各不相同，但如果把它们并排放，会发现每一处的失效条件都能翻译成同一句话：\emph{二阶项在你实际使用的尺度上不可忽略。}

Newton 法在拐点附近发散，因为分子 \(\abs{f'}\) 趋于零，可信半径塌到比当前误差还小，而算法照样按公式跨了一大步。梯度下降在曲率大的方向震荡，因为那个方向的可信半径最短，而步长是所有方向共用的一个数——震荡幅度由条件数决定，正是最长与最短半径之比。扩展 Kalman 滤波给出过度自信的协方差，因为协方差递推里用的是 Jacobian，它描述的是那条切平面的传播，真实映射的弯曲带来的额外不确定性根本没有进入账本；急转弯时弯曲最大，账本却最干净。delta 方法在样本量不足时失准，因为 \(\hat\theta\) 的波动范围超出了 \(g\) 的可信邻域；在 \(g'(\theta)=0\) 时它干脆给出零方差，那是分子塌到零的极端情形，正确做法是展到二阶。局部线性解释在决策边界附近不稳定，因为边界附近函数变化最剧烈，邻域取大一点结果就翻个个儿——同一个样本换个邻域半径得到相反的归因，这不是方法有 bug，是可信邻域比采样半径还窄。神经切核描述不了特征学习，因为特征学习恰恰意味着参数走出了线性化成立的范围。

这些不是八条经验，是一条数学事实的八次现身。所以调试时的问题也可以统一成一句：\emph{我实际迈的步子，有没有超出这一点的可信邻域。}多数时候答案是能查的——把步长减半再看一眼，行为若明显变好，就说明之前在邻域之外。

\begin{remark}
把可信半径写成 \(2\varepsilon\abs{f'}/\abs{f''}\) 是精确的读数，但用它来\emph{选}步长通常不现实：二阶导数本身要估，估它的代价常常不比直接试更低。实践中的做法是线搜索、信赖域、自适应步长，它们都不去算那个半径，而是通过``走一步，看看实际下降有没有达到线性模型的预测''来间接探测它。这一层是工程惯例而非定理，但它探测的对象仍然是同一个量。
\end{remark}

\subsection{什么时候线性化根本不该用}

有三类情形，减小步长也救不回来。函数在关心的点处不可微时，一阶项不存在——\(\relu\) 的折点、绝对值损失的零点、约束边界上的最优解都属于这一类，此时该换的工具是次梯度或近端算子而不是更小的步长。函数虽然光滑但曲率无界时，可信半径可以任意小，任何固定步长都会在某处失效；这解释了为什么强凸加上梯度 Lipschitz 这两条假设总是成对出现，它们一个给曲率下界一个给曲率上界，合起来把可信半径夹在一个正的区间里。第三类最隐蔽：被线性化的对象根本不是你以为的那个。用一个局部线性模型解释黑箱时，得到的是黑箱在这一点的一阶行为，不是它``为什么''这样判断；把前者当后者，减小邻域半径只会让结论更精确地回答一个不相干的问题。

线性化之所以值得单独讲一篇，是因为它给出的是同一个筹码的两面：一阶项让任意非线性的问题暂时变成线性问题，于是线性代数积累的全部工具立刻可用；而这份便利只在一个有限的邻域里成立，邻域的宽度还不由你决定。下一篇处理的是筹码的另一半——当问题\emph{已经}是线性的，它还能被简化到什么程度。答案是可以简化到几乎不剩下什么：换一组合适的坐标，一个纠缠的线性作用会变成一组互不干扰的独立缩放。
